%!%
% Copyright (c) 2026 mingcheng <mingcheng@outlook.com>
% 
% File: ch13-workflow.tex
% Author: mingcheng <mingcheng@outlook.com>
% File Created: 2026-03-27 15:50:02
% 
% Modified By: mingcheng <mingcheng@outlook.com>
% Last Modified: 2026-03-29 07:17:27
%%

% !TeX root = ../prompt-engineering-guide-for-beginners.tex

\chapter{从新手到高手：打造你的 AI 工作流}
\label{ch:workflow}

到这里，你已经掌握了提示工程的核心技巧。这一章要帮你从「会用 AI」升级到「善用 AI」，把零散的技巧整合成系统化的工作流。

\section{建立你的个人提示词库}

如果你用心写了一条效果很好的提示词，千万别用完就忘了。把它保存下来，下次遇到类似的任务直接复用。

\textbf{怎么建提示词库？}

不需要什么高级工具，一个文档或笔记本就够了。关键是做好分类和标注：

\begin{itemize}
  \item \textbf{分类标签}：按业务场景分类（如「邮件写作」「数据分析」「文案创作」），同时可以加上策略标签（如 \#零样本、\#链式思考），方便快速查找。
  \item \textbf{模板化设计}：不要把提示词写死，把需要替换的内容用方括号标出（如 \verb|[输入文本]|、\verb|[目标受众]|、\verb|[输出格式]|），下次使用时只需要填空就行。
  \item \textbf{记录效果与局限}：不仅记下好用的提示词，也记下它在什么情况下效果不好。这些经验比提示词本身更有价值。
  \item \textbf{版本记录}：像写代码一样记录每次优化的过程（如 v1 只有基本指令，v2 加入角色设定后效果更好），帮助你积累调优经验。
\end{itemize}

\begin{tipbox}
当提示词的积累达到一定量级时（超过 20 条高频词谱），请务必使用专门的知识管理工具如 Notion（含数据库视图）、Obsidian（含反向链接），甚至一些专门的 Github 提示词仓或第三方 Prompt SDK 服务，确保自己总能在一个统一搜索界面下，在 5 秒钟内复制并替换上相关的预设变量。
\end{tipbox}

\textbf{示例格式：}

\begin{promptbox}
\textbf{名称}：小红书种草文案

\textbf{场景}：为产品生成小红书风格的种草笔记

\textbf{提示词}：你是一位小红书爆款文案写手。请为以下产品写一篇种草笔记……（完整提示词）

\textbf{效果}：语气和格式都不错，偶尔需要手动调整 emoji 的使用密度

\textbf{版本}：v2（v1 没有限定字数，输出太长）
\end{promptbox}

积累到一定量之后，这个提示词库就是你的「AI 武器库」，能让你的工作效率实现质的飞跃。

\section{定制你的 AI 助手}

很多 AI 平台现在都提供了「自定义助手」的功能，让你可以创建一个预设好角色、知识和规则的专属 AI：

\begin{itemize}
  \item \textbf{ChatGPT 的 Custom GPTs}：你可以创建自己的 GPT，预设系统提示词、上传参考文档，甚至接入外部工具
  \item \textbf{Claude 的 Projects}：可以为特定项目创建独立的对话空间，上传相关文档作为知识库
  \item \textbf{Kimi 的智能体}：支持自定义角色和知识库，适合搭建中文场景下的专属助手
  \item \textbf{字节的 Coze}（扣子）：可以组合多种插件和工作流，搭建功能更复杂的 AI 应用
\end{itemize}

\textbf{定制 AI 助手的一些实用方向：}

\begin{itemize}
  \item \textbf{专属知识库问答}：上传你的行业研报、产品文档等资料，打造一位了解你业务的「专属分析师」。这种基于检索增强生成（RAG，参见第~\ref{sec:rag}~节）的方式能有效减少幻觉。
  \item \textbf{日常文书自动化}：打造一位帮你处理周报、会议纪要等固定格式文档的助手。通过预设好模板和示例，你只需输入关键信息就能得到规范的输出。
  \item \textbf{模拟演练}：创建一位扮演「刁钻客户」或「严格审稿人」的对抗角色，用来在正式提交前测试你的方案、文档是否经得起挑战。
  \item \textbf{工作流联动}：通过接入外部工具（如搜索、代码执行、消息推送等），让 AI 不仅能对话，还能实际执行操作，迈向智能体（Agent）架构。
\end{itemize}

\begin{tipbox}
定制 AI 助手的核心在于 System Prompt（系统提示词，参见第~\ref{sec:system-prompt}~节）。在这段设定指令中清晰定义助手的角色定位、回答流程和禁止行为，它的质量直接决定了助手的表现上限。
\end{tipbox}

\section{把 AI 融入你的工具链}

如果你目前使用 AI 的方式还停留在「复制 AI 的回答，粘贴到其他软件」，那还有很大的提效空间。更高效的做法是把 AI 当作工具链中的一个环节，和你日常使用的其他软件协作起来。

\textbf{常见的 AI + 工具组合：}

\begin{itemize}
  \item \textbf{AI + Excel/表格工具}：让 AI 帮你写复杂的 Excel 公式、生成数据分析的思路，或者直接处理你上传的表格数据。比如你不会写 VLOOKUP 公式，只需要告诉 AI「我有两张表，A 表有订单号和金额，B 表有订单号和客户名，请帮我写一个公式把客户名匹配到 A 表」，几秒钟就搞定了。
  \item \textbf{AI + 笔记工具}（Notion / 飞书文档 / Obsidian）：用 AI 快速生成初稿，再在笔记工具中整理和完善。很多笔记工具现在已经内置了 AI 功能，可以直接在编辑器里调用。
  \item \textbf{AI + 邮件}：用 AI 起草邮件，复制到邮箱客户端中微调后发送。特别适合需要用英文写邮件但又不太自信的场景。
  \item \textbf{AI + Python/代码}：让 AI 帮你写脚本来自动化重复性工作。即使你不会写代码也没关系，你只要用自然语言描述你要做的事情就行，比如「写一个脚本，把这个文件夹里所有 PDF 文件的文件名改成日期开头的格式」。
  \item \textbf{AI + 搜索引擎}：先用 AI 了解一个话题的基本框架和关键概念，再用搜索引擎查找具体数据和最新信息。AI 帮你建立认知框架，搜索引擎帮你验证和补充细节，两者配合效率最高。
  \item \textbf{AI + 设计工具}：先用 AI 生成文案和创意方向，再用 Canva、Figma 等设计工具落地视觉执行。你也可以让 AI 帮你写设计需求文档，然后把文档交给设计师或设计工具。
\end{itemize}

核心思路是：\textbf{让 AI 做它擅长的事情（生成初稿、处理文字、提供思路），让专业工具做它擅长的事情（精确计算、视觉设计、数据存储）}。不要试图让 AI 做所有事情，也不要忽略它能帮上忙的环节。两者结合，才能发挥最大价值。

\section{持续学习：跟上 AI 的进化速度}

AI 领域的变化速度非常快。新的模型、新的功能、新的技巧在不断涌现。今天的最佳实践，过几个月可能就有了更好的替代方案。

\textbf{如何保持学习：}

\begin{itemize}
  \item \textbf{定期尝试新工具}：每隔一段时间试试新出的 AI 工具或新功能。不需要深入研究，快速体验一下就能知道是否对你有用
  \item \textbf{关注几个可靠的信息源}：选择 2 到 3 个你信任的科技媒体或博主，定期浏览他们关于 AI 的更新
  \item \textbf{加入社区}：找一个 AI 爱好者的社区（微信群、论坛、Discord 等），和其他人交流使用心得
  \item \textbf{在工作流微调中实践（Do-by-Learning）}：不要为了学提示词技巧而硬刷。当你下一次在代码报错（Error Stack Trace）、大量网页去重整理（Data Processing）或是写重复信件时，暂停三秒反问：“能不能构建一段指令喂给模型批处理”，真实痛点中的磨砺胜过阅读万字干货。
  \item \textbf{无视参数量，聚焦逻辑拆解能力，关注大语言模型推理路线}：不必因没弄懂 O1 模型的底层原理（RLHF 机制）与最新参数量突破便感到自己落伍了！本质上，大模型在几年内依然是强大的“文字概率接龙机”，万金油法则便是锻炼自己将问题切分为极小的微元组件。
\end{itemize}

\begin{tipbox}
掌握结构化的提示词框架（就像写作中的「三段式」和编程中的「设计模式」）就足以应对 80\% 的工作挑战。至于模型底层的技术原理和 Token 优化等深入话题，留给专门研究这个方向的工程师就好。你的首要目标是让 AI 切实提升你的工作效率。
\end{tipbox}

记住，学习 AI 不是为了追赶潮流，而是为了让你的工作和生活更高效。找到适合自己的节奏，持续进步就好。
